\documentclass{article}
\usepackage{algorithm, algorithmic, amsfonts, amsmath}

\begin{document}
\begin{algorithm}[!h]
\renewcommand{\algorithmicrequire}{\textbf{Input:}} % show \REQUIRE as Input
\renewcommand{\algorithmicensure}{\textbf{Output:}} % show \ENSURE as Output
	\caption{TRP: future connection score on temporal graphs, trained by the nnPU risk (TPU)} % algorithm title
	\begin{algorithmic}[1] % [1] numbers every line
	    \REQUIRE
	    Temporal graph $G=\{G^{1},G^{2},\dots,G^{T}\}$ with node features $x_{v}^{t}$;
	    a node pair $a_{i,j}=\langle v_{i},v_{j}\rangle$ in $G^{t}$;
	    an initialized pair-embedding vector $h_{a_{i,j}}^{0}$ (e.g. zeros)
	    \ENSURE
	    $p_{a_{i,j}}^{t}$, the connectivity prediction score for the node pair $a_{i,j}$
	    \FOR{$\tau\leftarrow 1\cdots t$}
	    \STATE Obtain the current node features $x_{v}^{\tau}$ ($v=v_{i},v_{j}$), as well as
	    $x_{Nr(v)}^{\tau}$ for the sampled neighbouring nodes of $v_{i},v_{j}$
	    \STATE Aggregate the neighbourhood information:
	    $z_{v}^{\tau}=f_{G}\!\left(x_{v}^{\tau},x_{Nr(v)}^{\tau};\theta_{G}\right)$,\quad $v=v_{i},v_{j}$ % this paper aggregates by max pooling
	    \STATE Update the pair embedding:
	    $h_{a_{i,j}}^{\tau}=f_{A}\!\left(h_{a_{i,j}}^{\tau-1},z_{v_{i}}^{\tau},z_{v_{j}}^{\tau};\theta_{A}\right)$ % recursion over time, capturing how the pairwise relation evolves
	    \ENDFOR
	    \RETURN $p_{a_{i,j}}^{t}=f_{C}\!\left(h_{a_{i,j}}^{t};\theta_{C}\right)$
	    \\\hrulefill
	    \STATE \textbf{Training:} fit the score above with the \emph{non-negative PU risk}
	    (Kiryo et al.), because a currently unconnected pair is unlabeled, not negative % an unconnected pair now does not mean it will stay unconnected
	    \STATE \textbf{Testing:} run the same procedure with $t=T$ to obtain $p_{a_{i,j}}^{T}$
	\end{algorithmic}
\end{algorithm}
% Akujuobi, Chen, Elhoseiny, Spranger, Zhang, NeurIPS 2020, "Temporal Positive-unlabeled
% Learning for Biomedical Hypothesis Generation via Risk Estimation".
% Compiled from Algorithm 1 of the original; the training objective and test procedure are in its Section 3.
% Naming: the original calls the temporal representation module TRP and the whole method TPU. Figure 2 of the survey
% uses TRP, so the directory name follows the figure.
% Task: biomedical hypothesis generation, predicting whether two scientific terms will become associated in future.
% Two PU points:
% (1) a term pair unobserved now is not a negative, only not yet studied, so a non-negative risk estimator is used
%     instead of ordinary binary classification; (2) relations evolve, so the model updates pair embeddings
%     recursively over tau = 1..t rather than reading a single static snapshot.
% Difference from others in the branch: GPL addresses heterophily bias in node classification, PULL the PU structure
% of edge prediction, and TRP temporal evolution; the three target different graph PU problems.
\end{document}
